Диссертация посвящена послойному и вычислительно сопоставимому исследованию
режимов предиктивного кодирования (PC) относительно обратного распространения
ошибки (BP). Работа разделяет четыре аналитических измерения, которые часто
смешиваются при сравнении алгоритмов обучения: сходство конечного поведения,
сходство внутреннего механизма, вычислительную стоимость и возможность
до штатного завершения заменить оставшуюся каноническую итеративную часть
зарегистрированным ограниченным аналитическим завершением при сохранении
требуемого ответа. Для перехода
от наблюдения внутреннего состояния к вычислительному решению вводятся две
ограниченные теоретические конструкции: PC-TREF задаёт эквивалентность
состояний относительно требуемого действия, полную достаточность
диагностического представления на зарегистрированной области и более узкую
селективную достаточность правила
для конкретного действия, а PC-CATM описывает формирование, перенос и
компенсацию каналов коррекции. Архитектура QWake-PC операционализирует эти основания; в работе
экспериментально проверяется её ограниченная реализация QWake-FP для FixedPred.
В QWake-FP раннее действие означает выбор зарегистрированного аналитического
завершения \nolinkurl{fixedpred_eta1_wavefront_completion_v1} вместо
оставшихся канонических итеративных проходов; полный точный остаток
\nolinkurl{complete_suffix_stage2_baseline_v1} служит эталонным и
резервным путём.

Экспериментальная программа организована как последовательность заранее
зафиксированных этапов. Этапы этап~1 и этап~2 сравнивают BP, Exact, FixedPred
и Strict в контролируемых условиях. этап~3A исследует послойные градиенты и
представления; этап~3B локализует стоимость итеративного вывода состояния и
проверяет кандидаты точного вычислительного замещения посредством заранее
зарегистрированных проверок эквивалентности. Финальная ветка QWake-FP разделяет
три последовательных вопроса: присутствуют ли на поверхности траекторий с
зафиксированной целостностью предтерминальные записи, состояния которых допускают
раннее действие относительно задачи; можно ли на зафиксированной калибровочной
поверхности селективно распознавать ненулевое подмножество таких записей по состоянию,
доступному до действия, без наблюдавшихся опасных принятий; и даёт ли такое
решение положительную выгоду при
заранее зафиксированном полном учёте стоимости принятия решения.

В исследованных условиях FixedPred демонстрировал более высокую наблюдаемую
близость к BP по направлению градиента и внутренним представлениям, чем Strict,
при одновременном уменьшении нормы
градиента в ранних слоях. Профилирование этап~3B связало существенную часть
стоимости исследованных режимов PC с итеративным выводом состояния; последующие
кандидаты B1/B2 прошли зарегистрированные проверки эквивалентности, а полная
матрица сопоставленного профилирования была завершена без повторных попыток.
При этом зарегистрированный инженерный экран продолжения дал 0/16
квалифицированных конфигураций во всех четырёх группах B1/B2×FixedPred/Strict,
и оба кандидата получили решение \texttt{reject\_or\_revise}; этот результат не
интерпретируется как утверждение о превосходстве исходной линии. В QWake-FP поверхность с зафиксированной целостностью содержала предтерминальные
записи, состояния которых эталонная проверка относила к классу
\texttt{EARLY\_ADMISSIBLE}, а ненулевая селективная распознаваемость их
подмножества по состоянию до действия с нулём наблюдавшихся опасных принятий
наблюдалась внутри замороженного семейства из 2625 скалярных правил принятия
решения на 756 калибровочных записях. Из них 264 правила имели ненулевое
покрытие при отсутствии опасных принятых решений, а правило с наибольшим
покрытием среди правил с нулём наблюдавшихся опасных принятий покрывало 216
из 756 калибровочных записей (28.57\%). Само правило использует только
\texttt{compute\_step >= 5}; поэтому положительный результат C08 на этой
поверхности задаёт временную границу, эквивалентную фиксированному префиксу, а
не демонстрирует зависимость решения от содержания входа или преимущество
механизмных признаков PC-CATM. Структурная сверка временной поверхности
показывает, что эти 216 записей состоят из 108 предтерминальных записей шага 5 с
одним оставшимся проходом и 108 записей терминальной границы шага 6; поэтому
28.57\% сохраняется как зарегистрированное покрытие полной поверхности и не
трактуется как доля предтерминальных применений раннего аналитического действия. При этом ни одно правило
не обеспечило положительную совокупную чистую экономию при зафиксированном
полном учёте стоимости принятия решения. В полной стоимости этого правила
преобладали накладные расходы наблюдателя; правило C2 не было зафиксировано,
поэтому подтверждающий этап C3 не открывался.

Полученный отрицательный экономический результат относится именно к полному
учёту стоимости принятия решения, зафиксированному для C2, и не измеряет
добавочную стоимость выполнения минимального распознавателя. Поэтому работа
поддерживает вывод о наличии в исследованной области предтерминальных записей,
состояния которых допускают раннее действие относительно задачи, и ненулевой
селективной калибровочной распознаваемости их подмножества по состоянию до
действия с нулём наблюдавшихся опасных принятий;
отклоняет гипотезу о существовании
правила с нулём наблюдавшихся опасных принятий и положительной экономической
выгодой в замороженном семействе C2 при
указанном учёте стоимости; и оставляет экономическую осуществимость минимальной
реализации отдельной непроверенной оцениваемой величиной.

\noindent\textbf{Ключевые слова:} предиктивное кодирование, обратное
распространение ошибки, послойный анализ, вычислительная стоимость,
аналитическое завершение, трассируемость.
